\documentclass[12pt,a4paper]{article}
\usepackage[T1]{fontenc}
\usepackage[utf8]{inputenc}
\usepackage[french]{babel}
\usepackage{geometry}
\geometry{margin=2cm}
\usepackage{tikz}
\usetikzlibrary{shapes,positioning,arrows.meta,calc,fit,backgrounds}
\usepackage{caption}
\usepackage{amsmath}
\usepackage{enumitem}

\tikzset{
  entity/.style={
    rectangle,draw,thick,fill=blue!10,
    minimum width=2.8cm,minimum height=0.6cm,
    align=center,font=\small\bfseries,
    inner sep=2pt
  },
  attr/.style={
    rectangle,draw,thin,fill=blue!3,
    minimum width=2.8cm,align=left,font=\footnotesize,
    inner sep=3pt,anchor=north
  },
  assoc/.style={
    diamond,draw,thick,fill=orange!15,
    minimum width=2cm,minimum height=1cm,
    align=center,font=\footnotesize\bfseries,
    inner sep=1pt,aspect=2
  },
  note/.style={
    rectangle,draw,dashed,thin,fill=yellow!10,
    minimum width=2cm,align=center,font=\scriptsize\itshape,
    inner sep=3pt
  },
  card/.style={font=\scriptsize,fill=white,inner sep=1pt},
  rel/.style={->,>=stealth,thick},
  relb/.style={<-,>=stealth,thick},
  bi_rel/.style={<->,>=stealth,thick},
  dashed_rel/.style={->,>=stealth,thick,dashed}
}

\title{\textbf{Modèle Conceptuel de Données (MCD)\\[0.3em]
\large Application SDK WOOD Mobile}}
\author{}
\date{}

\begin{document}
\maketitle
\thispagestyle{empty}

\section{Introduction}

Ce document présente le Modèle Conceptuel de Données (MCD) de l'application SDK WOOD Mobile, construit selon la méthode Merise. Le MCD identifie les entités, leurs attributs (avec clés primaires soulignées et clés étrangères annotées), les associations et les cardinalités.

Les tables \texttt{vue\_depots} et \texttt{vue\_produits} sont des vues SQL utilisées comme clés étrangères mais ne constituent pas des entités du MCD ; elles sont représentées en pointillé.

Le MCD est organisé en cinq modules métier correspondant aux sprints de développement.

\tableofcontents
\newpage

% ================================================================
\section{Module 1 : Authentification \& Permissions}
% ================================================================

\subsection{Entités}

\begin{itemize}[leftmargin=1.5em]
\item \textbf{UTILISATEUR} (\underline{id}, nom, prenom, role, ValidationRetourChauffeur)
\item \textbf{UTILISATEUR\_TOKEN} (\underline{id}, idUser $\rightarrow$ UTILISATEUR, token)
\item \textbf{MENUS} (\underline{id}, code, designation)
\item \textbf{PERMISSIONS} (\underline{id}, user\_id $\rightarrow$ UTILISATEUR, menu\_id $\rightarrow$ MENUS, can\_view, can\_create, can\_update, can\_delete, can\_take\_picture\_bl, can\_close\_bl, can\_achever)
\end{itemize}

\subsection{Associations}

\begin{itemize}[leftmargin=1.5em]
\item \textbf{Possède\_token} : UTILISATEUR $1,n$ — $0,n$ UTILISATEUR\_TOKEN
\item \textbf{A\_permissions} : UTILISATEUR $0,n$ — $0,n$ MENUS (via PERMISSIONS, porte 7 attributs booléens)
\end{itemize}

\subsection{Diagramme}

\begin{figure}[ht]
\centering
\begin{tikzpicture}[node distance=2.2cm]

% --- Entities ---
\node[entity] (user) at (0,0) {UTILISATEUR};
\node[attr, below=0.15cm of user] (user_a) {\textbf{\underline{id}}, nom, prenom,\\role, ValidationRetourChauffeur};

\node[entity] (token) at (6,0.5) {TOKEN};
\node[attr, below=0.15cm of token] (token_a) {\textbf{\underline{id}}, idUser $\to$, token};

\node[entity] (menus) at (0,-3.5) {MENUS};
\node[attr, below=0.15cm of menus] (menus_a) {\textbf{\underline{id}}, code, designation};

\node[entity] (perms) at (5.5,-3.5) {PERMISSIONS};
\node[attr, below=0.15cm of perms] (perms_a) {\textbf{\underline{id}}, user\_id $\to$, menu\_id $\to$,\\can\_view, can\_create, can\_update,\\can\_delete, can\_take\_picture\_bl,\\can\_close\_bl, can\_achever};

% --- Associations ---
\node[assoc] (tok_assoc) at (3,0.3) {Possède};
\node[assoc] (perm_assoc) at (2.7,-3.5) {A\_permissions};

% --- Edges ---
\draw[bi_rel] (user) -- node[card,above]{$1,n$} (tok_assoc);
\draw[bi_rel] (tok_assoc) -- node[card,above]{$0,n$} (token);

\draw[bi_rel] (user) -- node[card,right]{$0,n$} (perm_assoc);
\draw[bi_rel] (perm_assoc) -- node[card,above]{$0,n$} (perms);

\end{tikzpicture}
\caption{MCD — Module Authentification \& Permissions}
\label{fig:mcd-auth}
\end{figure}

\newpage

% ================================================================
\section{Module 2 : Voyages \& Bons de Livraison}
% ================================================================

\subsection{Entités}

\begin{itemize}[leftmargin=1.5em]
\item \textbf{MARQUE\_VEHICULE} (\underline{id}, designation)
\item \textbf{VEHICULE} (\underline{id}, immatriculation, idMarque $\rightarrow$ MARQUE\_VEHICULE)
\item \textbf{VILLE} (\underline{id}, designation, kmDeCasa)
\item \textbf{VOYAGE} (\underline{id}, date\_depart, km\_depart, km\_retour, date\_retour, date\_create, statut, idChauffeur $\rightarrow$ UTILISATEUR, idVehicule $\rightarrow$ VEHICULE, idCreate $\rightarrow$ UTILISATEUR, depot\_depart $\rightarrow$ \textit{vue\_depots}, idVille $\rightarrow$ VILLE)
\item \textbf{VOYAGE\_BL} (\underline{id}, statut, coordinates, idVoyage $\rightarrow$ VOYAGE, idBL $\rightarrow$ DOCUMENT)
\item \textbf{VOYAGE\_BL\_IMAGE} (\underline{id}, nom\_fichier, chemin\_fichier, date\_upload, idVoyageBL $\rightarrow$ VOYAGE\_BL)
\item \textbf{DOCUMENT} (\underline{id}, id\_document, id\_type\_document, datetime\_document, id\_entreprise $\rightarrow$ PARTENAIRES)
\item \textbf{DOCUMENT\_ITEMS} (\underline{id\_item}, id\_document $\rightarrow$ DOCUMENT, idproduit $\rightarrow$ \textit{vue\_produits})
\item \textbf{PARTENAIRES} (\underline{id}, societe)
\end{itemize}

\subsection{Notes}
\begin{itemize}[leftmargin=1.5em]
\item \texttt{vue\_depots} (vue SQL) : identifié par \texttt{id}, attribut \texttt{nom}. Référencé par \texttt{depot\_depart} et \texttt{idDepotSource/Destination}.
\item \texttt{vue\_produits} (vue SQL) : identifié par \texttt{produit\_id}, attributs \texttt{produit}, \texttt{unite\_v}, \texttt{type\_stock}.
\item \texttt{document.id\_type\_document = 7} correspond aux Bons de Livraison (BL).
\end{itemize}

\subsection{Diagramme}

\begin{figure}[ht]
\centering
\begin{tikzpicture}[node distance=1.8cm]

% --- Entities ---
\node[entity] (user) at (0,2) {UTILISATEUR};

\node[entity] (voyage) at (0,-1.5) {VOYAGE};
\node[attr, below=0.1cm of voyage] (voyage_a) {\underline{id}, date\_depart, km\_depart,\\km\_retour, date\_retour, statut,\\idChauffeur $\to$, idVehicule $\to$,\\idCreate $\to$, depot\_depart $\to$, idVille $\to$};

\node[entity] (vehicule) at (-4.5,-1.5) {VEHICULE};
\node[attr, below=0.1cm of vehicule] (veh_a) {\underline{id}, immatriculation,\\idMarque $\to$};

\node[entity] (marque) at (-4.5,-5) {MARQUE};
\node[attr, below=0.1cm of marque] (marque_a) {\underline{id}, designation};

\node[entity] (ville) at (5,-1.5) {VILLE};
\node[attr, below=0.1cm of ville] (ville_a) {\underline{id}, designation, kmDeCasa};

\node[entity] (vbl) at (0,-6.5) {VOYAGE\_BL};
\node[attr, below=0.1cm of vbl] (vbl_a) {\underline{id}, statut, coordinates,\\idVoyage $\to$, idBL $\to$};

\node[entity] (vblimg) at (-3.5,-6.5) {VBL\_IMAGE};
\node[attr, below=0.1cm of vblimg] (vblimg_a) {\underline{id}, nom\_fichier,\\chemin\_fichier, date\_upload,\\idVoyageBL $\to$};

\node[entity] (doc) at (4,-6.5) {DOCUMENT};
\node[attr, below=0.1cm of doc] (doc_a) {\underline{id}, id\_document,\\id\_type\_document,\\datetime\_document,\\id\_entreprise $\to$};

\node[entity] (part) at (8,-6.5) {PARTENAIRES};
\node[attr, below=0.1cm of part] (part_a) {\underline{id}, societe};

\node[note] (depot_note) at (-2.5,2.6) {[vue\_depots :\\id, nom]};

% --- Associations ---
\node[assoc] (a_chauffeur) at (1.5,0.3) {Effectue};
\node[assoc] (a_cree) at (-1.5,0.3) {Crée};
\node[assoc] (a_marque) at (-4.5,-3.2) {De};
\node[assoc] (a_ville) at (3,-1.5) {Destination};
\node[assoc] (a_depot) at (-1.2,2.2) {Départ};
\node[assoc] (a_contient) at (1.8,-4) {Contient};
\node[assoc] (a_photo) at (-1.8,-5.3) {A\_photos};
\node[assoc] (a_entreprise) at (6,-6.5) {Entreprise};

% --- Edges ---
\draw[bi_rel] (user) -- node[card,right]{$0,n$} (a_chauffeur);
\draw[bi_rel] (a_chauffeur) -- node[card,right]{$1,n$} (voyage);

\draw[bi_rel] (user) -- node[card,left]{$0,n$} (a_cree);
\draw[bi_rel] (a_cree) -- node[card,left]{$0,n$} (voyage);

\draw[bi_rel] (vehicule) -- node[card,left]{$n,1$} (a_marque);
\draw[bi_rel] (a_marque) -- node[card,left]{$1,n$} (marque);

\draw[bi_rel] (voyage) -- node[card,above]{$0,1$} (a_ville);
\draw[bi_rel] (a_ville) -- node[card,above]{$0,n$} (ville);

\draw[dashed_rel] (voyage) -- node[card,above left]{$1,1$} (a_depot);
\draw[dashed_rel] (a_depot) -- (depot_note);

\draw[bi_rel] (voyage.south) -- node[card,left]{$1,n$} (a_contient.north);
\draw[bi_rel] (a_contient.south) -- node[card,right]{$0,n$} (vbl.north);

\draw[bi_rel] (vbl) -- node[card,above]{$1,n$} (a_photo);
\draw[bi_rel] (a_photo) -- node[card,above]{$0,n$} (vblimg);

\draw[bi_rel] (vbl) -- node[card,above]{$0,n$} (doc);
\draw[bi_rel] (doc) -- node[card,above]{$1,1$} (part);

\end{tikzpicture}
\caption{MCD — Module Voyages \& Bons de Livraison}
\label{fig:mcd-voyage}
\end{figure}

\newpage

% ================================================================
\section{Module 3 : Retours Chauffeur \& Rapports Qualité}
% ================================================================

\subsection{Entités}

\begin{itemize}[leftmargin=1.5em]
\item \textbf{RETOUR\_CHAUFFEUR} (\underline{id}, Bl\_cachetet, reglement, retour\_Mse, reclamation, statut, date, commentaire, client\_id $\rightarrow$ PARTENAIRES, chauffeur\_id $\rightarrow$ UTILISATEUR)
\item \textbf{RETOUR\_CHAUFFEUR\_IMAGE} (\underline{id}, nom\_fichier, chemin\_fichier, date\_upload, idCreate $\rightarrow$ UTILISATEUR, idRetourChaufeur $\rightarrow$ RETOUR\_CHAUFFEUR)
\item \textbf{RAPPORT\_QUALITE} (\underline{id}, dum, dossier, commentaire, dateCreate, dateModif, user\_id $\rightarrow$ UTILISATEUR, idCreate $\rightarrow$ UTILISATEUR, idModif $\rightarrow$ UTILISATEUR)
\item \textbf{RAPPORT\_QUALITE\_IMAGE} (\underline{id}, nom\_fichier, chemin\_fichier, date\_upload, idCreate $\rightarrow$ UTILISATEUR, idRapportQualite $\rightarrow$ RAPPORT\_QUALITE)
\item \textbf{PARTENAIRES} (\underline{id}, societe)
\item \textbf{UTILISATEUR} (\underline{id}, nom, prenom, role, ValidationRetourChauffeur)
\end{itemize}

\subsection{Diagramme}

\begin{figure}[ht]
\centering
\begin{tikzpicture}[node distance=2cm]

% --- Entities ---
\node[entity] (user) at (0,1) {UTILISATEUR};

\node[entity] (rc) at (-4,-2.5) {RETOUR\_CHAUFFEUR};
\node[attr, below=0.1cm of rc] (rc_a) {\underline{id}, Bl\_cachetet, reglement,\\retour\_Mse, reclamation, statut,\\date, commentaire,\\client\_id $\to$, chauffeur\_id $\to$};

\node[entity] (rc_img) at (-7.5,-2.5) {RC\_IMAGE};
\node[attr, below=0.1cm of rc_img] (rc_img_a) {\underline{id}, nom\_fichier,\\chemin\_fichier, date\_upload,\\idCreate $\to$, idRetourChaufeur $\to$};

\node[entity] (rq) at (4,-2.5) {RAPPORT\_QUALITE};
\node[attr, below=0.1cm of rq] (rq_a) {\underline{id}, dum, dossier,\\commentaire, dateCreate, dateModif,\\user\_id $\to$, idCreate $\to$, idModif $\to$};

\node[entity] (rq_img) at (8,-2.5) {RQ\_IMAGE};
\node[attr, below=0.1cm of rq_img] (rq_img_a) {\underline{id}, nom\_fichier,\\chemin\_fichier, date\_upload,\\idCreate $\to$, idRapportQualite $\to$};

\node[entity] (part) at (-4,-6.5) {PARTENAIRES};
\node[attr, below=0.1cm of part] (part_a) {\underline{id}, societe};

% --- Associations ---
\node[assoc] (a_soumet) at (-2.5,-0.8) {Soumet};
\node[assoc] (a_photo_rc) at (-5.8,-2.5) {A\_photos};
\node[assoc] (a_create_prc) at (-7.5,0.1) {Photo\_de};
\node[assoc] (a_client) at (-4,-4.5) {Client};
\node[assoc] (a_rapport_user) at (2,-0.8) {Concerne};
\node[assoc] (a_rapport_create) at (4,0.1) {Crée};
\node[assoc] (a_rapport_modif) at (5.5,0.8) {Modifie};
\node[assoc] (a_photo_rq) at (6,-2.5) {A\_photos};
\node[assoc] (a_create_prq) at (8,0.1) {Photo\_de};

% --- Edges Retour ---
\draw[bi_rel] (user) -- node[card,above]{$0,n$} (a_soumet);
\draw[bi_rel] (a_soumet) -- node[card,below]{$0,n$} (rc);

\draw[bi_rel] (rc) -- node[card,above]{$1,n$} (a_photo_rc);
\draw[bi_rel] (a_photo_rc) -- node[card,above]{$0,n$} (rc_img);

\draw[bi_rel] (user) -- node[card,left]{$0,n$} (a_create_prc);
\draw[bi_rel] (a_create_prc) -- node[card,left]{$0,n$} (rc_img);

\draw[bi_rel] (rc) -- node[card,left]{$0,n$} (a_client);
\draw[bi_rel] (a_client) -- node[card,left]{$1,1$} (part);

% --- Edges Rapport ---
\draw[bi_rel] (user) -- node[card,above]{$0,n$} (a_rapport_user);
\draw[bi_rel] (a_rapport_user) -- node[card,below]{$0,n$} (rq);

\draw[bi_rel] (user) -- node[card,above right]{$0,n$} (a_rapport_create);
\draw[bi_rel] (a_rapport_create) -- node[card,right]{$0,n$} (rq);

\draw[bi_rel] (user) -- node[card,right]{$0,n$} (a_rapport_modif);
\draw[bi_rel] (a_rapport_modif) -- node[card,right]{$0,n$} (rq);

\draw[bi_rel] (rq) -- node[card,above]{$1,n$} (a_photo_rq);
\draw[bi_rel] (a_photo_rq) -- node[card,above]{$0,n$} (rq_img);

\draw[bi_rel] (user) -- node[card,above]{$0,n$} (a_create_prq);
\draw[bi_rel] (a_create_prq) -- node[card,below]{$0,n$} (rq_img);

\end{tikzpicture}
\caption{MCD — Module Retours Chauffeur \& Rapports Qualité}
\label{fig:mcd-retour}
\end{figure}

\newpage

% ================================================================
\section{Module 4 : Lieux de Projets}
% ================================================================

\subsection{Entités}

\begin{itemize}[leftmargin=1.5em]
\item \textbf{PROJET} (\underline{id}, projet, localisation, commentaire, contact\_nom, contact\_telephone, idCreate $\rightarrow$ UTILISATEUR)
\item \textbf{PROJET\_IMAGE} (\underline{id}, nom\_fichier, chemin\_fichier, date\_upload, idProjet $\rightarrow$ PROJET)
\item \textbf{UTILISATEUR} (\underline{id}, nom, prenom, role)
\end{itemize}

\subsection{Diagramme}

\begin{figure}[ht]
\centering
\begin{tikzpicture}[node distance=2.5cm]

\node[entity] (user) at (0,0) {UTILISATEUR};

\node[entity] (projet) at (0,-3) {PROJET};
\node[attr, below=0.1cm of projet] (projet_a) {\underline{id}, projet, localisation,\\commentaire, contact\_nom,\\contact\_telephone, idCreate $\to$};

\node[entity] (pimg) at (5,-3) {PROJET\_IMAGE};
\node[attr, below=0.1cm of pimg] (pimg_a) {\underline{id}, nom\_fichier,\\chemin\_fichier, date\_upload,\\idProjet $\to$};

\node[assoc] (a_create) at (0,-1.5) {Crée};
\node[assoc] (a_photo) at (2.5,-3) {A\_photos};

\draw[bi_rel] (user) -- node[card,right]{$0,n$} (a_create);
\draw[bi_rel] (a_create) -- node[card,right]{$0,n$} (projet);

\draw[bi_rel] (projet) -- node[card,above]{$1,n$} (a_photo);
\draw[bi_rel] (a_photo) -- node[card,above]{$0,n$} (pimg);

\end{tikzpicture}
\caption{MCD — Module Lieux de Projets}
\label{fig:mcd-projet}
\end{figure}

\newpage

% ================================================================
\section{Module 5 : Demande de Transfert}
% ================================================================

\subsection{Entités}

\begin{itemize}[leftmargin=1.5em]
\item \textbf{DEMANDETRANSFERT} (\underline{id}, reference, dum, transporteur, matricule, date, observation, statut, dateSingeDemandeur, dateEnvoie, dateRecu, idDepotSource $\rightarrow$ \textit{vue\_depots}, idDepotDestination $\rightarrow$ \textit{vue\_depots}, idCreate $\rightarrow$ UTILISATEUR, idSigneDemamdeur $\rightarrow$ UTILISATEUR, idEnvoie $\rightarrow$ UTILISATEUR, idRecu $\rightarrow$ UTILISATEUR)
\item \textbf{DEMANDETRANSFERT\_DETAILS} (\underline{id}, nbrFDX, Qte, Unite, preparer, idDT $\rightarrow$ DEMANDETRANSFERT, idProduit $\rightarrow$ \textit{vue\_produits})
\item \textbf{TRANSFER\_LOTS} (\underline{idItem}, \underline{idProduit}, \underline{Lot}, preparer, qte, old\_lot, idModif $\rightarrow$ UTILISATEUR, idItem $\rightarrow$ DEMANDETRANSFERT\_DETAILS)
\item \textbf{UTILISATEUR} (\underline{id}, nom, prenom, role)
\end{itemize}

\subsection{Notes}
\begin{itemize}[leftmargin=1.5em]
\item \texttt{vue\_depots} : référencé par \texttt{idDepotSource} et \texttt{idDepotDestination} (dépôt source et destination du transfert).
\item \texttt{vue\_produits} : référencé par \texttt{idProduit} dans \texttt{DEMANDETRANSFERT\_DETAILS} et \texttt{TRANSFER\_LOTS}.
\item \texttt{TRANSFER\_LOTS} a une clé primaire composée : (\underline{idItem}, \underline{idProduit}, \underline{Lot}).
\item \texttt{DEMANDETRANSFERT} a 4 rôles distincts vers \texttt{UTILISATEUR} : Crée, Signe, Envoie, Reçoit.
\end{itemize}

\subsection{Diagramme}

\begin{figure}[ht]
\centering
\begin{tikzpicture}[node distance=2cm]

% --- Entities ---
\node[entity] (user) at (0,1.5) {UTILISATEUR};

\node[entity] (dt) at (0,-2) {DEMANDE\_TRANSFERT};
\node[attr, below=0.1cm of dt] (dt_a) {\underline{id}, reference, dum, transporteur,\\matricule, date, observation, statut,\\dateSingeDemandeur, dateEnvoie, dateRecu,\\idDepotSource $\to$, idDepotDestination $\to$,\\idCreate $\to$, idSigneDemamdeur $\to$,\\idEnvoie $\to$, idRecu $\to$};

\node[entity] (dd) at (5,-2) {DT\_DETAILS};
\node[attr, below=0.1cm of dd] (dd_a) {\underline{id}, nbrFDX, Qte, Unite,\\preparer, idDT $\to$, idProduit $\to$};

\node[entity] (tl) at (5,-6) {TRANSFER\_LOTS};
\node[attr, below=0.1cm of tl] (tl_a) {\underline{idItem}, \underline{idProduit}, \underline{Lot},\\preparer, qte, old\_lot,\\idModif $\to$};

\node[note] (depot_note) at (-3.5,1.8) {[vue\_depots :\\id, nom]};
\node[note] (prod_note) at (9,-2) {[vue\_produits :\\produit\_id, produit,\\unite\_v, type\_stock]};

% --- Associations ---
\node[assoc] (a_create) at (-1.2,0.2) {Crée};
\node[assoc] (a_signe) at (0.5,0.8) {Signe};
\node[assoc] (a_envoie) at (-1.5,0.9) {Envoie};
\node[assoc] (a_recu) at (1,0.2) {Reçoit};
\node[assoc] (a_depot_s) at (-2.5,-0.2) {Dépôt\_S};
\node[assoc] (a_depot_d) at (2.5,-0.2) {Dépôt\_D};
\node[assoc] (a_detail) at (2.5,-2) {Contient};
\node[assoc] (a_lot) at (5,-4) {A\_lots};
\node[assoc] (a_modif_lot) at (7,-4.5) {Modifie\_lot};

% --- Edges ---
\draw[bi_rel] (user) -- node[card,left]{$0,n$} (a_create);
\draw[bi_rel] (a_create) -- node[card,left]{$0,n$} (dt);

\draw[bi_rel] (user) -- node[card,right]{$0,n$} (a_signe);
\draw[bi_rel] (a_signe) -- node[card,right]{$0,n$} (dt);

\draw[bi_rel] (user) -- node[card,left]{$0,n$} (a_envoie);
\draw[bi_rel] (a_envoie) -- node[card,below]{$0,n$} (dt);

\draw[bi_rel] (user) -- node[card,right]{$0,n$} (a_recu);
\draw[bi_rel] (a_recu) -- node[card,right]{$0,n$} (dt);

\draw[dashed_rel] (dt) -- node[card,above left]{$0,n$} (a_depot_s);
\draw[dashed_rel] (a_depot_s) -- (depot_note);

\draw[dashed_rel] (dt) -- node[card,above right]{$0,n$} (a_depot_d);
\draw[dashed_rel] (a_depot_d) -- (depot_note);

\draw[bi_rel] (dt) -- node[card,above]{$1,n$} (a_detail);
\draw[bi_rel] (a_detail) -- node[card,above]{$0,n$} (dd);

\draw[dashed_rel] (dd) -- node[card,above]{$0,n$} (prod_note);

\draw[bi_rel] (dd) -- node[card,left]{$1,n$} (a_lot);
\draw[bi_rel] (a_lot) -- node[card,left]{$0,n$} (tl);

\draw[bi_rel] (user) -- node[card,above right]{$0,n$} (a_modif_lot);
\draw[bi_rel] (a_modif_lot) -- node[card,below]{$0,n$} (tl);

\end{tikzpicture}
\caption{MCD — Module Demande de Transfert}
\label{fig:mcd-transfert}
\end{figure}

\newpage

% ================================================================
\section{Dictionnaire de données récapitulatif}
% ================================================================

\subsection{Entités et leurs attributs}

\begin{center}
\renewcommand{\arraystretch}{1.2}
\begin{tabular}{|p{3.2cm}|p{2cm}|p{8cm}|}
\hline
\textbf{Entité} & \textbf{Clé primaire} & \textbf{Attributs} \\
\hline
UTILISATEUR & \underline{id} & nom, prenom, role, ValidationRetourChauffeur \\
\hline
UTILISATEUR\_TOKEN & \underline{id} & idUser $\to$, token \\
\hline
MENUS & \underline{id} & code, designation \\
\hline
PERMISSIONS & \underline{id} & user\_id $\to$, menu\_id $\to$, can\_view, can\_create, can\_update, can\_delete, can\_take\_picture\_bl, can\_close\_bl, can\_achever \\
\hline
MARQUE\_VEHICULE & \underline{id} & designation \\
\hline
VEHICULE & \underline{id} & immatriculation, idMarque $\to$ \\
\hline
VILLE & \underline{id} & designation, kmDeCasa \\
\hline
VOYAGE & \underline{id} & date\_depart, km\_depart, km\_retour, date\_retour, date\_create, statut, idChauffeur $\to$, idVehicule $\to$, idCreate $\to$, depot\_depart $\to$, idVille $\to$ \\
\hline
VOYAGE\_BL & \underline{id} & statut, coordinates, idVoyage $\to$, idBL $\to$ \\
\hline
VOYAGE\_BL\_IMAGE & \underline{id} & nom\_fichier, chemin\_fichier, date\_upload, idVoyageBL $\to$ \\
\hline
RETOUR\_CHAUFFEUR & \underline{id} & Bl\_cachetet, reglement, retour\_Mse, reclamation, statut, date, commentaire, client\_id $\to$, chauffeur\_id $\to$ \\
\hline
RETOUR\_CHAUFFEUR\_IMAGE & \underline{id} & nom\_fichier, chemin\_fichier, date\_upload, idCreate $\to$, idRetourChaufeur $\to$ \\
\hline
RAPPORT\_QUALITE & \underline{id} & dum, dossier, commentaire, dateCreate, dateModif, user\_id $\to$, idCreate $\to$, idModif $\to$ \\
\hline
RAPPORT\_QUALITE\_IMAGE & \underline{id} & nom\_fichier, chemin\_fichier, date\_upload, idCreate $\to$, idRapportQualite $\to$ \\
\hline
PROJET & \underline{id} & projet, localisation, commentaire, contact\_nom, contact\_telephone, idCreate $\to$ \\
\hline
PROJET\_IMAGE & \underline{id} & nom\_fichier, chemin\_fichier, date\_upload, idProjet $\to$ \\
\hline
DEMANDETRANSFERT & \underline{id} & reference, dum, transporteur, matricule, date, observation, statut, dateSingeDemandeur, dateEnvoie, dateRecu, idDepotSource $\to$, idDepotDestination $\to$, idCreate $\to$, idSigneDemamdeur $\to$, idEnvoie $\to$, idRecu $\to$ \\
\hline
DEMANDETRANSFERT\_DETAILS & \underline{id} & nbrFDX, Qte, Unite, preparer, idDT $\to$, idProduit $\to$ \\
\hline
TRANSFER\_LOTS & \underline{idItem}, \underline{idProduit}, \underline{Lot} & preparer, qte, old\_lot, idModif $\to$ \\
\hline
DOCUMENT & \underline{id} & id\_document, id\_type\_document, datetime\_document, id\_entreprise $\to$ \\
\hline
DOCUMENT\_ITEMS & \underline{id\_item} & id\_document $\to$, idproduit $\to$ \\
\hline
PARTENAIRES & \underline{id} & societe \\
\hline
\end{tabular}
\end{center}

\subsection{Règles de gestion}

\begin{enumerate}[leftmargin=1.5em]
\item Un \textbf{utilisateur} possede de 0 à $n$ tokens d'authentification.
\item Un \textbf{utilisateur} a de 0 à $n$ permissions, chacune associée à un \textbf{menu}.
\item Un \textbf{voyage} est effectué par un et un seul \textbf{chauffeur} (utilisateur).
\item Un \textbf{voyage} utilise un et un seul \textbf{véhicule}, d'une \textbf{marque}.
\item Un \textbf{voyage} part d'un \textbf{dépôt} et peut avoir une \textbf{ville} de destination.
\item Un \textbf{voyage} contient de 0 à $n$ \textbf{bons de livraison} (documents de type 7).
\item Chaque \textbf{BL de voyage} peut avoir de 0 à $n$ \textbf{photos}.
\item Une \textbf{retour chauffeur} est soumis par un \textbf{chauffeur} et concerne un \textbf{client} (partenaire).
\item Un \textbf{retour chauffeur} a de 0 à $n$ \textbf{images}, créées par un \textbf{utilisateur}.
\item Un \textbf{rapport qualité} est soumis, créé et modifié par des \textbf{utilisateurs} (3 rôles distincts).
\item Un \textbf{rapport qualité} a de 0 à $n$ \textbf{images}, créées par un \textbf{utilisateur}.
\item Un \textbf{projet} est créé par un \textbf{utilisateur} et a de 0 à $n$ \textbf{images}.
\item Une \textbf{demande de transfert} est créée, signée, envoyée et recue par des \textbf{utilisateurs} (4 rôles distincts).
\item Une \textbf{demande de transfert} contient de 0 à $n$ \textbf{détails} (lignes de produits).
\item Chaque \textbf{détail} peut avoir de 0 à $n$ \textbf{lots}.
\item Un \textbf{document} appartient à une \textbf{entreprise} (partenaire) et contient des \textbf{articles}.
\item Les vues SQL \texttt{vue\_depots} et \texttt{vue\_produits} sont référencées par des FK mais ne constituent pas des entités du MCD.
\end{enumerate}

\end{document}